\chapter{PEG Tube Placement}
\index{PEG Tube Placement}

\medskip
\noindent\textit{This chapter is for general education. Follow the individual instructions supplied by the nutrition, nursing, and medical team.}

\section*{What it is}
Percutaneous endoscopic gastrostomy (PEG) placement creates an opening through the abdominal wall into the stomach so prescribed liquid nutrition, fluids, or medicines can be given. It may be recommended when swallowing is unsafe or oral intake is not enough. It does not automatically mean that eating or drinking by mouth must stop.

\section*{Before the procedure}
The team should discuss the reason, alternatives, expected benefits, consent, anaesthesia, bleeding and infection risks, and what support will be available at home. Tell them about allergies, medicines that affect bleeding, diabetes, pregnancy possibility, and previous abdominal surgery. Fasting and medicine instructions must be followed.

\section*{After placement}
A nurse or dietitian should teach the person and carers how to position for feeding, flush the tube, give medicines safely, clean the skin, check the external marking, and obtain supplies. Keep the head and upper body raised during feeding and for the period advised by the team. Do not change the feed, water volume, or tube position without guidance.

\section*{Common problems}
Mild soreness or a small amount of clear discharge can occur while the site heals. Infection, leakage, blockage, displacement, bleeding, aspiration, and buried-bumper syndrome are possible. Never force a blocked tube or give a feed if position is uncertain.

\section*{When to Seek Medical Care}
Stop feeding and seek urgent advice for pain during feeding or flushing, new bleeding, significant leakage, repeated vomiting, breathing difficulty, a displaced or fallen-out tube, severe abdominal pain, fever, or marked redness and swelling around the site.

\section*{Sources and further reading}
\begin{itemize}
\item Leeds Teaching Hospitals NHS Trust. \textit{PEG feeding tube care advice}. \url{https://www.leedsth.nhs.uk/patients/resources/peg-feeding-tube-care-advice/}
\item University College London Hospitals. \textit{Gastrostomy or PEG feeding}. \url{https://www.uclh.nhs.uk/patients-and-visitors/patient-information-pages/gastrostomy-peg-feeding-making-decision}
\end{itemize}
